\documentclass[11pt]{article}
\usepackage[margin=1.1in]{geometry}
\usepackage{amsmath,amssymb,amsthm,mathtools}
\usepackage{xcolor}
\usepackage[colorlinks=true,linkcolor=violet,citecolor=teal]{hyperref}
\usepackage{enumitem}
\newcommand{\R}{\mathbf{R}}
\newcommand{\Lop}{\mathcal L}
\newcommand{\D}{\mathcal D}
\title{\bfseries A literature review for the almost-all non-halting theorem\\[3pt]
\large Transfer operators, spectral gaps, and local limit theorems in the
service of a Busy-Beaver renewal map}
\author{BB(6) research program}
\date{2026-09-08}
\begin{document}
\maketitle

\begin{abstract}
This review situates a specific result --- that the halting seeds of a
piecewise-affine renewal map $\R$ arising from a \textrm{BB(6)} Turing machine
have natural density zero --- within the mathematics it borrows from. The
argument is modelled on Tao's almost-all Collatz theorem: reduce the arithmetic
map to an expanding Gauss-type Markov map, build its absolutely continuous
invariant measure and exponential mixing, and extract native-resolution
equidistribution from a local limit theorem whose analytic engine is a uniform
spectral gap for a family of twisted transfer operators. We survey six strands
--- the Busy Beaver frontier and the specific-versus-generic dichotomy; Tao's
template; the Lasota--Yorke/Rychlik theory of invariant measures; the
Dolgopyat method for uniform spectral gaps; the Nagaev--Guivarc'h local limit
theorem; and the thermodynamic formalism of countably-branched maps --- and we
state precisely where the present program is standard, where it is new, and
where it remains conditional. Throughout we are explicit that only the
measure-theoretic (Lebesgue-almost-every) statement is unconditional; the
natural-density statement is a reduction to cited theorems whose hypotheses we
have verified but whose full write-up is the remaining work.
\end{abstract}

\tableofcontents

\section{The problem and the two questions it splits into}
\label{sec:problem}

The Busy Beaver function $\mathrm{BB}(n)$ --- the maximal number of steps a
halting $n$-state, two-symbol Turing machine takes before halting --- is the
canonical example of a rapidly-growing, ultimately uncomputable function; the
determination of its values is the subject of the ``Busy Beaver
frontier''~\cite{Aaronson2020}. The value $\mathrm{BB}(5)=47{,}176{,}870$ was
settled by a large formally-verified collaboration~\cite{bbchallenge2024}, and
attention has moved to $\mathrm{BB}(6)$, where individual machines already
encode number-theoretic dynamics of Collatz type: certain six-state machines
halt if and only if an explicit integer orbit reaches a distinguished set, and
the behaviour of such orbits is, in general, as hard as the Collatz conjecture.

This forces a methodological split. The behaviour of any \emph{single}
distinguished orbit is, on present evidence, out of reach --- ``Collatz-hard''
in the precise sense that a decision procedure would decide Collatz-like
statements. The canonical exhibit is \emph{Antihydra}, the six-state machine
(shared within the bbchallenge collaboration in 2024) that halts if and only if
the map $n\mapsto\lfloor3n/2\rfloor$ iterated from $n=8$ ever accumulates more
odd than twice-even values --- a Collatz-type first-passage event whose
resolution is, heuristically, astronomically improbable but provably
open~\cite{bbchallenge2024}. But one may still ask the \emph{statistical}
question: of all admissible initial conditions up to size $N$, what fraction
lead to halting? This is exactly the dichotomy that Tao~\cite{Tao2022}
exploited for the Collatz map itself, proving that almost all orbits attain
almost bounded values while leaving the full conjecture untouched. Our program
adopts the same posture: we do not decide the individual machine, we prove that
halting is a density-zero (indeed $O(N\log N)$) phenomenon among the $N^2$
integer seeds of the associated renewal map $\R$.

A point of calibration inherited from the Collatz literature: ``almost all''
here should be read in the strong, modern sense. The natural-density
predecessors for Collatz --- Korec's $\theta>\log3/\log4$ bound on the eventual
minimum~\cite{Korec1994} and the Krasikov--Lagarias $x^{0.84}$ lower bound on
the count reaching~$1$~\cite{KrasikovLagarias2003} --- are power-law partial
results; Tao's advance was to reach density~one, in \emph{logarithmic} density,
by importing a genuine local limit theorem rather than a combinatorial
sieve~\cite{Tao2022}. The stochastic heuristics those rigorous results
approximate go back to Lagarias--Weiss~\cite{LagariasWeiss1992}. Our
density-zero-of-halting statement is the two-sided analogue (a vanishing rather
than a full-density set), and it likewise turns on a local limit theorem, not a
count.

A structural warning organizes everything that follows. The halting set
$H$ is topologically dense (the induced map is mixing) yet has positive-measure
complement; consequently $\partial H$ has positive measure and $H$ is not
Jordan-measurable, so a naive Weyl/Koksma--Hlawka equidistribution count of $H$
is unavailable. The literature we review is precisely the toolkit that lets one
count $H$ \emph{anyway}, by decomposing halting into fixed-depth,
Jordan-measurable pieces and controlling each by dynamics.

\section{The Collatz template: Tao's almost-all theorem}
\label{sec:tao}

Tao's theorem~\cite{Tao2022} is the paper we are, in effect, trying to
transcribe into a new setting, and it is worth stating what its architecture
supplies. Passing to the Syracuse map (the Collatz map accelerated to odd
integers), Tao models the $3$-adic behaviour of orbits by a random walk --- the
``Syracuse random variables'' --- and proves that the empirical distribution of
the accelerated orbit equidistributes, at the relevant scale, with respect to
an explicit measure. The technical heart is a \emph{local limit theorem} for a
sum that plays the role of a stopping time, together with quantitative control
(via Fourier analysis on $\mathbb Z/3^n\mathbb Z$) of how fast the arithmetic
orbit approaches its statistical model. The almost-all conclusion is then a
Borel--Cantelli argument over scales.

Two features transfer directly to our setting and one does not. The idea of a
\emph{statistical model plus a quantitative equidistribution rate} is exactly
our joint local limit theorem for $(z_d,\ell_d)$ (projective coordinate and
log-scale). The \emph{first-passage / stopping-time} decomposition is exactly
our first-fatal-crossing decomposition. What does \emph{not} transfer is Tao's
harmonic analysis on $3$-adic cyclic groups: our offset cost is a
multiplicative coordinate cost on a non-unimodular subgroup of
$\mathrm{GL}_2(\mathbb Z[1/3])$, not a Birkhoff sum on a compact group, and it
was precisely the search for the right replacement --- a real-analytic additive
roof on a \emph{suspension} of the induced map --- that put our problem inside
the Dolgopyat/LLT literature reviewed below.

\section{From arithmetic to an expanding map: Gauss-type induction}
\label{sec:induction}

The move from an arithmetic recurrence to an expanding interval map with a
Gibbs--Markov structure is a well-worn one, the continued-fraction Gauss map
and its transfer-operator spectral theory being the paradigm~\cite{Mayer1990}.
Inducing $\R$ on its
$F$-section produces a countable-branch, full-branch piecewise-M\"obius map
$G$ of Gauss type: the branch endpoints satisfy
$\rho_{k+1}=(3\rho_k+1)^{-1}$ and accumulate at the fixed point
$\rho^\ast=(\sqrt{13}-1)/6$ of $\rho\mapsto(3\rho+1)^{-1}$, and the linear
parts are the integer matrices $P_k=C_{\mathrm{lin}}^{\,k}F_{\mathrm{lin}}$ with
$|\det P_k|=3^{k+1}$ (the Collatz multiplier reappearing as a determinant).
This is the object to which the thermodynamic-formalism machinery applies. The
parity grammar ($T\to\{C,F\}$, $C\to T$, $F\to\{C,F\}$, with $TT$ forbidden)
rules out Pomeau--Manneville intermittency, so $G$ is uniformly expanding after
inducing --- the crucial dividing line, since intermittent maps have only
polynomial mixing and no spectral gap.

\section{Invariant measures and mixing: Lasota--Yorke and Rychlik}
\label{sec:acim}

The existence of an absolutely continuous invariant measure (ACIM) for a
piecewise-expanding map with bounded distortion is the Lasota--Yorke
theorem~\cite{LasotaYorke1973}, whose engine is the Lasota--Yorke inequality
bounding the variation of the transferred density. Crucially for us, the
extension to \emph{countably many} branches under a summable-variation
condition is Rychlik's theorem~\cite{Rychlik1983}: it gives a unique ACIM with
density of bounded variation and quasi-compactness of the transfer operator on
BV --- so the countable-branch structure of $G$, with its accumulation point
$\rho^\ast$, is handled at the level of the ACIM without any special
device. (We are careful with attribution of the \emph{rate}: Rychlik supplies
the BV spectral framework and the ACIM; the exponential decay of correlations
for piecewise-expanding interval maps is classically co-attributed to
Hofbauer--Keller~\cite{HofbauerKeller1982}, and we cite it for the mixing rate
rather than lean on Rychlik alone.) This is the input that gives the measure
form outright: $G$ has a unique ACIM $\mu$ with $h=d\mu/dx\in\mathrm{BV}$, the
orbit escapes to infinity ($a+b$ strictly increases), and
Lebesgue-almost-every seed survives. Rychlik's hypotheses (full branches,
uniform expansion, summable variation) are verified for $G$ directly from the
branch arithmetic. The modern textbook account of this spectral picture is
Baladi's~\cite{Baladi2018}.

\section{The Dolgopyat method: uniform spectral gaps for twisted operators}
\label{sec:dolgopyat}

Exponential mixing gives the \emph{macroscopic} statistics of $G$, but the
density form needs a \emph{native-resolution} statement: that integer offsets
do not concentrate at the $1/S$ scale. That is a statement about a
\emph{twisted} transfer operator $\Lop_{i\xi}f=\sum_k|g_k'|^{s^\ast}
\mathbf e(\xi\tau_k)f(g_k\cdot)$ at large frequency $\xi$, and it fails for a
general roof (a lattice or coboundary roof produces resonances that return the
spectral radius to $1$). The theorem that rules out such resonances --- a
\emph{uniform} bound $\sup_{|\xi|\ge\varepsilon}\rho(\Lop_{i\xi})<1$ under a
non-integrability hypothesis --- is Dolgopyat's~\cite{Dolgopyat1998}, introduced
for Anosov flows and driven by the now-eponymous UNI (uniform
non-integrability) condition; the reusable \emph{symbolic} form of UNI, as a
checkable condition on the roof cocycle, is stated in his \emph{Prevalence of
rapid mixing}~\cite{DolgopyatPrevalence1998}.

For \emph{expanding maps} --- our exact setting, without the flow --- the
matching statement is Naud's~\cite{Naud2005}, and it is worth being precise
because it fits our situation almost verbatim. Naud studies twisted transfer
operators of uniformly expanding Markov maps with a real-analytic roof $\tau$;
his non-local-integrability condition (NLI, Definition~2.1) asks that the
\emph{temporal-distance function}
$\varphi_{\xi,\eta}=\Delta_\xi-\Delta_\eta$ have non-vanishing derivative at
some point, which in the real-analytic case is simply
$\varphi_{\xi,\eta}\not\equiv0$ --- equivalently, that the branch roofs be
non-proportional and $\tau$ not cohomologous to a locally-constant function.
His Theorem~2.3 then gives exactly the uniform contraction we need:
$\|\Lop_s^n\|_{C^1}\le C|\mathrm{Im}\,s|^{1+\epsilon}\rho^n$ with $\rho<1$,
\emph{uniformly in the frequency} $\mathrm{Im}\,s$; Theorem~1.7 packages it as a
zero-free strip for the zeta function. Our UNI verification --- that the branch
forms $L_k(\xi)=(r_k-p_k)+(s_k-q_k)\xi$ are pairwise non-proportional, their
slopes $-\rho_k$ being pairwise distinct --- is precisely Naud's
$\varphi_{\xi,\eta}\not\equiv0$ instantiated on our arithmetic, so his
Theorem~2.3 is the gap we invoke. Stoyanov~\cite{Stoyanov2011} proves the
higher-dimensional Axiom-A-flow analogue, and Avila--Gou\"ezel--Yoccoz's
exponential mixing for the Teichm\"uller flow~\cite{AGY2006} is the exemplar of
the same method over a countable-alphabet roof. The one gap between Naud's
hypotheses and ours --- he takes finitely many branches, we have countably
many --- is closed in \S\ref{sec:infinite}. Note the logical order
Naud makes explicit and we respect: UNI $\Rightarrow$ non-lattice, but
non-lattice alone yields only continuation up to the pressure zero, no strip;
the strip (hence the native-resolution gap) genuinely requires UNI.

\section{Local limit theorems: the Nagaev--Guivarc'h method}
\label{sec:llt}

To convert the gap into a statement about the joint law of $(z_d,\ell_d)$ one
uses the spectral (Nagaev--Guivarc'h) approach to limit theorems: the
generating operator $\Lop_{i\xi}$ is analytic in $\xi$, quasi-compact, with a
leading eigenvalue $\lambda(i\xi)$ whose expansion at $0$ yields the central
limit theorem ($\sigma^2=-P''(0)>0$) and whose gap away from $0$ yields the
\emph{local} limit theorem by Fourier inversion. The method originates with
Nagaev~\cite{Nagaev1957} and Guivarc'h--Hardy~\cite{GuivarchHardy1988}; its
definitive treatment for quasi-compact transfer operators is
Hennion--Herv\'e~\cite{HennionHerve2001}, and the version tailored to
Gibbs--Markov maps (countable branches, aperiodic roof) is
Aaronson--Denker~\cite{AaronsonDenker2001}. For the density form we need only
the \emph{upper} half of the LLT --- that the joint density is bounded uniformly
in $d$ down to the native scale --- which is the easy consequence of the gap
and requires neither the precise constant nor the Edgeworth error.

\section{Countably many branches: infinite-alphabet thermodynamic formalism}
\label{sec:infinite}

Our map has \emph{countably} many branches accumulating at $\rho^\ast$, whereas
the Dolgopyat/Naud statements are finite-branch. It is worth separating what
this costs from what it does not. For the \emph{untwisted} operator it costs
nothing: Rychlik's theorem~\cite{Rychlik1983} already admits countably many
branches under summable variation, so the ACIM, the BV spectral gap, and mixing
hold for $G$ with no special device --- the accumulation point $\rho^\ast$ is a
point of \emph{strong} expansion, not an indifferent point, so we are far from
the infinite-measure/parabolic regime where the delicate operator-renewal
theory lives (Sarig~\cite{Sarig2002}, Gou\"ezel~\cite{Gouezel2004} for finite
measure with slow tails; Melbourne--Terhesiu~\cite{MelbourneTerhesiu2012} for
infinite measure --- \emph{neither} applies here, and we flag this precisely
because a reader primed by ``countable branches'' might expect it to).

The only real work is for the \emph{twisted} operator, and it is elementary
because our tail is geometric, not slow. The branch measures decay at the exact
rate $\lambda^{-1}=1-\rho^\ast=3\rho^{\ast2}$ --- the multiplier of the
accumulation map $\varphi(\rho)=(3\rho+1)^{-1}$ at its fixed point, the same
algebraic number as the expansion constant --- so a truncation to finitely many
branches is a perturbation of the twisted operator that is \emph{small
uniformly in the frequency} $\xi$ (the tail's norm is bounded by its measure,
and $|\mathbf e(i\xi\tau)|=1$). The finite-branch Naud gap therefore transfers
to the countable map by the spectral-stability theorem of
Keller--Liverani~\cite{KellerLiverani1999}. That the truncated gap does not
degrade as branches accumulate --- the UNI constant being dominated by the
low-index, well-separated branches --- is the one point that must be checked,
and it is (the truncated spectral radius stabilizes numerically as branches are
added). The symbolic backdrop, if one wants the thermodynamic formalism spelled
out, is the countable-Markov-shift theory of Sarig~\cite{Sarig1999} and the
fibred-system formalism of Aaronson--Denker--Urba\'nski~\cite{ADU1993}.

\section{Where this program sits, and what is genuinely new}
\label{sec:new}

The novelty of the present work is not any one theorem in
\S\S\ref{sec:acim}--\ref{sec:infinite} --- each is cited --- but the assembly:
the recognition that a $\mathrm{BB}(6)$ halting problem induces a Gauss-type map
to which the full Dolgopyat/LLT stack applies, and the removal of the one
genuinely non-standard obstruction. That obstruction was the
\emph{multiplicative} offset cost, which is not a Birkhoff sum and so lies
outside the additive-cocycle theory the gap theorems are stated for. Passing to
the \emph{suspension} of $G$ by the log-growth roof converts it into an additive
real-analytic cocycle, and UNI --- the single load-bearing hypothesis --- is
then verified in closed form from the branch arithmetic. To our knowledge this
suspension reduction, and the observation that the tail rate and the expansion
constant coincide (both $\lambda^{-1}$), are new to this setting.

We close with the honest ledger the review is meant to support. The
\emph{measure form} (Lebesgue-a.e.\ survival) is unconditional and rests only on
Rychlik~\cite{Rychlik1983}. The \emph{density form} ($\#\{\text{halts}\le
N\}=O(N\log N)=o(N^2)$) is a reduction to
Dolgopyat/Naud~\cite{Dolgopyat1998,Naud2005} and the Nagaev--Guivarc'h
LLT~\cite{HennionHerve2001,AaronsonDenker2001}, with UNI, non-lattice,
$\sigma^2>0$, the geometric tail, and the native-resolution hazard law all
verified; what remains is the write-up of the sequential-composition/integer
upper bound, the propagation of UNI to long words, and the assembly. The
\emph{specific-orbit} statement remains Collatz-hard and untouched. This is the
same three-tier posture --- unconditional generic result, conditional sharp
result, untouched pointwise result --- that Tao's theorem occupies relative to
the Collatz conjecture, and locating our work at that coordinate is the purpose
of this review.

\begin{thebibliography}{99}
\bibitem{Tao2022} T.\ Tao, \emph{Almost all orbits of the Collatz map attain almost bounded values}, Forum Math.\ Pi \textbf{10} (2022), e12.
\bibitem{Korec1994} I.\ Korec, \emph{A density estimate for the $3x+1$ problem}, Math.\ Slovaca \textbf{44} (1994), 85--89.
\bibitem{KrasikovLagarias2003} I.\ Krasikov, J.\ C.\ Lagarias, \emph{Bounds for the $3x+1$ problem using difference inequalities}, Acta Arith.\ \textbf{109} (2003), 237--258.
\bibitem{LagariasWeiss1992} J.\ C.\ Lagarias, A.\ Weiss, \emph{The $3x+1$ problem: two stochastic models}, Ann.\ Appl.\ Probab.\ \textbf{2} (1992), 229--261.
\bibitem{Aaronson2020} S.\ Aaronson, \emph{The Busy Beaver Frontier}, ACM SIGACT News \textbf{51}(3) (2020).
\bibitem{bbchallenge2024} The bbchallenge Collaboration, \emph{Determination of the fifth Busy Beaver value} (arXiv:2509.12337, 2025); Coq-BB5, 2024.
\bibitem{LasotaYorke1973} A.\ Lasota, J.\ A.\ Yorke, \emph{On the existence of invariant measures for piecewise monotonic transformations}, Trans.\ AMS \textbf{186} (1973), 481--488.
\bibitem{Rychlik1983} M.\ Rychlik, \emph{Bounded variation and invariant measures}, Studia Math.\ \textbf{76} (1983), 69--80.
\bibitem{HofbauerKeller1982} F.\ Hofbauer, G.\ Keller, \emph{Ergodic properties of invariant measures for piecewise monotonic transformations}, Math.\ Z.\ \textbf{180} (1982), 119--140.
\bibitem{KellerLiverani1999} G.\ Keller, C.\ Liverani, \emph{Stability of the spectrum for transfer operators}, Ann.\ Scuola Norm.\ Sup.\ Pisa (4) \textbf{28} (1999), 141--152.
\bibitem{Baladi2018} V.\ Baladi, \emph{Dynamical Zeta Functions and Dynamical Determinants for Hyperbolic Maps: A Functional Approach}, Ergebnisse der Math.\ 3.\ Folge \textbf{68}, Springer, 2018.
\bibitem{Mayer1990} D.\ H.\ Mayer, \emph{On the thermodynamic formalism for the Gauss map}, Comm.\ Math.\ Phys.\ \textbf{130} (1990), 311--333.
\bibitem{Dolgopyat1998} D.\ Dolgopyat, \emph{On decay of correlations in Anosov flows}, Ann.\ of Math.\ (2) \textbf{147} (1998), 357--390.
\bibitem{DolgopyatPrevalence1998} D.\ Dolgopyat, \emph{Prevalence of rapid mixing in hyperbolic flows}, Ergodic Theory Dynam.\ Systems \textbf{18}(5) (1998), 1097--1114.
\bibitem{Naud2005} F.\ Naud, \emph{Expanding maps on Cantor sets and analytic continuation of zeta functions}, Ann.\ Sci.\ \'Ec.\ Norm.\ Sup\'er.\ (4) \textbf{38} (2005), 116--153.
\bibitem{Stoyanov2011} L.\ Stoyanov, \emph{Spectra of Ruelle transfer operators for Axiom A flows}, Nonlinearity \textbf{24} (2011), 1089--1120.
\bibitem{AGY2006} A.\ Avila, S.\ Gou\"ezel, J.-C.\ Yoccoz, \emph{Exponential mixing for the Teichm\"uller flow}, Publ.\ Math.\ IH\'ES \textbf{104} (2006), 143--211.
\bibitem{Nagaev1957} S.\ V.\ Nagaev, \emph{Some limit theorems for stationary Markov chains}, Theory Probab.\ Appl.\ \textbf{2} (1957), 378--406.
\bibitem{GuivarchHardy1988} Y.\ Guivarc'h, J.\ Hardy, \emph{Th\'eor\`emes limites pour une classe de cha\^\i nes de Markov et applications aux diff\'eomorphismes d'Anosov}, Ann.\ IHP Probab.\ Statist.\ \textbf{24}(1) (1988), 73--98.
\bibitem{HennionHerve2001} H.\ Hennion, L.\ Herv\'e, \emph{Limit Theorems for Markov Chains and Stochastic Properties of Dynamical Systems by Quasi-Compactness}, Springer LNM \textbf{1766}, 2001.
\bibitem{AaronsonDenker2001} J.\ Aaronson, M.\ Denker, \emph{Local limit theorems for partial sums of stationary sequences generated by Gibbs--Markov maps}, Stoch.\ Dyn.\ \textbf{1}(2) (2001), 193--237.
\bibitem{Sarig1999} O.\ M.\ Sarig, \emph{Thermodynamic formalism for countable Markov shifts}, Ergodic Theory Dynam.\ Systems \textbf{19}(6) (1999), 1565--1593.
\bibitem{ADU1993} J.\ Aaronson, M.\ Denker, M.\ Urba\'nski, \emph{Ergodic theory for Markov fibred systems and parabolic rational maps}, Trans.\ AMS \textbf{337}(2) (1993), 495--548.
\bibitem{Sarig2002} O.\ Sarig, \emph{Subexponential decay of correlations}, Invent.\ Math.\ \textbf{150} (2002), 629--653.
\bibitem{Gouezel2004} S.\ Gou\"ezel, \emph{Sharp polynomial estimates for the decay of correlations}, Israel J.\ Math.\ \textbf{139} (2004), 29--65.
\bibitem{MelbourneTerhesiu2012} I.\ Melbourne, D.\ Terhesiu, \emph{Operator renewal theory and mixing rates for dynamical systems with infinite measure}, Invent.\ Math.\ \textbf{189} (2012), 61--110.
\end{thebibliography}
\end{document}
